\chapter{اقتصاد گذار و مدیریت دارایی‌های ملی}
\label{app:economy}

مدیریت اقتصادی در دوران گذار باید بین سه اولویت توازن برقرار کند: تثبیت اضطراری، بازسازی اعتماد و توسعه بلندمدت.

\section{برنامه‌ی اقتصادی سه‌مرحله‌ای}

\begin{modernbox}{فاز ۱: تثبیت اضطراری (۱۰۰ روز اول)}
هدف: جلوگیری از فروپسانی معیشت و تأمین کالاهای اساسی.
\begin{itemize}
    \item حسابرسی فوری بانک مرکزی و تثبیت نرخ ارز.
    \item تضمین زنجیره تأمین غذا و دارو.
    \item پرداخت یارانه‌های نقدی اضطراری به دهک‌های پایین.
    \item مذاکره با نهادهای بین‌المللی (IMF و بانک جهانی) برای بسته‌های حمایتی.
\end{itemize}
\end{modernbox}

\begin{modernbox}{فاز ۲: بازسازی و اصلاحات (ماه ۳ تا ۱۸)}
\begin{itemize}
    \item بازسازی شفاف بودجه ملی و حذف هزینه‌های ایدئولوژیک.
    \item اصلاح نظام مالیاتی و بانکی.
    \item جذب سرمایه‌گذاری خارجی و فراخوان به دیاسپورای ایرانی.
    \item خصوصی‌سازی شفاف نهادهای وابسته به سپاه و بنیادها.
\end{itemize}
\end{modernbox}

\section{فرایند بازیابی دارایی‌های غارت‌شده}

تخمین زده می‌شود صدها میلیارد دلار از ثروت ملی در اختیار نهادهای غیرپاسخ‌گو و افراد وابسته به نظام باشد.

\begin{center}
\begin{tikzpicture}[node distance=3cm, auto]
    \tikzstyle{block} = [rectangle, draw, fill=SoftRed!20, text width=3cm, text centered, rounded corners, minimum height=3em, font=\small]
    \tikzstyle{line} = [draw, -Latex, line width=1pt]

    \node [block] (audit) {\textbf{۱. حسابرسی} \\ شناسایی اموال در داخل و خارج};
    \node [block, left of=audit] (freeze) {\textbf{۲. مسدودسازی} \\ توقیف دارایی‌های مشکوک};
    \node [block, below of=freeze] (recover) {\textbf{۳. بازیابی} \\ فرایندهای حقوقی بین‌المللی};
    \node [block, right of=recover] (fund) {\textbf{۴. تخصیص} \\ واریز به صندوق ملی بازسازی};

    \path [line] (audit) -- (freeze);
    \path [line] (freeze) -- (recover);
    \path [line] (recover) -- (fund);
\end{tikzpicture}
\end{center}

\section{تخصیص منابع در صندوق ملی بازسازی}

\begin{center}
\begin{tikzpicture}
    \pie[rotate=0, radius=2.5, text=inside, color={SoftPurple, SoftGreen, LightBlue, SoftGold, SlateGray}]
    {30/جبران خسارت قربانیان, 30/پروژه‌های زیرساختی, 20/صندوق نسل‌های آینده, 10/آموزش و بهداشت, 10/محیط زیست}
\end{tikzpicture}
\end{center}

\begin{center}
\begin{longtable}{R{4cm} R{4cm} R{6cm}}
\caption{وضعیت اقتصادی در لحظه‌ی گذار} \\
\hline
\textbf{چالش‌ها} & \textbf{دارایی‌ها} & \textbf{اقدام فوری} \\
\hline
\endfirsthead
\hline
\textbf{چالش‌ها} & \textbf{دارایی‌ها} & \textbf{اقدام فوری} \\
\hline
\endhead
\hline
\endfoot
تورم شدید و سقوط ریال & نیروی کار تحصیل‌کرده & کنترل نرخ ارز و واردات کالا \\
فساد سیستماتیک & منابع عظیم انرژی و معدن & حسابرسی بنیادها و نهادها \\
بیکاری گسترده & موقعیت ژئوپلیتیک تجاری & پروژه‌های عمرانی کوچک و سریع \\
تحریم‌های بین‌المللی & دیاسپورای ثروتمند و متخصص & دیپلماسی فعال برای لغو تحریم‌ها \\
\hline
\end{longtable}
\end{center}
